\documentclass[11pt]{article}

\usepackage[a4paper,margin=1in]{geometry}
\usepackage[T1]{fontenc}
\usepackage[utf8]{inputenc}
\usepackage{lmodern}
\usepackage{amsmath,amssymb,amsthm}
\usepackage[hidelinks]{hyperref}
\usepackage{booktabs}
\usepackage{xcolor}

\title{Set-Period vs.\ Multiset-Period: Computational Results}
\author{Dirk Kunert}
\date{\today}

\begin{document}

\maketitle

\section*{Setup}

For each row $(o_n, o_d, a_n, a_d, \mathrm{period})$ of the file
\texttt{tests/new\_find\_patterns\_x\_max\_1000000\_51012\_lines.csv}, the
program \texttt{src/c/set/add\_period\_set} computed the set-based period
length $\lambda_{\text{set}}$ on the same parameter quadruple and appended
it as a sixth column~\texttt{period\_set}. The run was performed in
\emph{global} mode with $X_{\max}=10^{6}$ and a per-row wall-clock cap of
$120$~seconds. The output is written to
\texttt{tests/new\_find\_patterns\_x\_max\_1000000\_51012\_lines.csv.tmp}.

\section*{Aggregate counts}

Of the $51{,}011$ data rows in the input, the resulting file breaks down
as follows.

\begin{center}
\begin{tabular}{lrl}
\toprule
Outcome & Rows & Meaning \\
\midrule
$\lambda_{\text{set}} = \mathrm{period}$  & $50{,}968$ & set period equals multiset period \\
$\lambda_{\text{set}} = -2$ & $29$ & \texttt{ARRAY\_SIZE\_EXCEEDED} \\
$\lambda_{\text{set}} = -3$ & $1$  & \texttt{DX\_LENGTH\_TO\_SMALL} \\
$\lambda_{\text{set}} = -4$ & $13$ & per-row $120$\,s timeout \\
$\lambda_{\text{set}} > 0$, $\lambda_{\text{set}} \neq \mathrm{period}$ & $1$ & genuine disagreement (see below) \\
\midrule
Total                                     & $51{,}012$ & \\
\bottomrule
\end{tabular}
\end{center}

In other words, on $99.92\%$ of the rows ($50{,}968 / 51{,}011$) the
set-based and multiset-based period lengths coincide; the remaining
$43$ rows split into $43$ computational aborts (out of memory, dx-window
too small, or timeout) and exactly one row in which the two periods are
both well-defined positive integers but disagree.

\section*{The single positive disagreement}

The only row with $\lambda_{\text{set}} > 0$ and
$\lambda_{\text{set}} \neq \mathrm{period}$ is

\begin{center}
\begin{tabular}{rrrrrr}
\toprule
$o_n$ & $o_d$ & $a_n$ & $a_d$ & multiset period & $\lambda_{\text{set}}$ \\
\midrule
$98{,}281$ & $139$ & $68$ & $149$ & $153{,}431$ & $1$ \\
\bottomrule
\end{tabular}
\end{center}

Here the multiset period is $153{,}431$ while the set-based period
collapses to $1$. This is a candidate witness for the
multiset-vs.-set distinction discussed in Open Problem~2 and warrants a
direct re-computation by hand or with an independent implementation
before being cited.

\section*{Run summary}

\begin{itemize}
\item Resume from a partial run: $40{,}201$ rows already present in the
      output were skipped; $10{,}811$ rows were freshly processed.
\item No hard failures during the resume run; $11$ of the $13$ timeouts
      shown above were produced in this run, the remaining $2$ were
      already present in the partial output.
\item Errors of types $-2$ and $-3$ accumulated $30$ rows in total;
      these are computational limitations of the current
      \texttt{lambda} implementation at $X_{\max}=10^{6}$ and do not
      indicate a true mismatch of the two period notions.
\end{itemize}

\section*{Caveats}

\begin{itemize}
\item Negative \texttt{period\_set} values mean the set-based period was
      \emph{not} computed, not that it differs from the multiset period.
      Excluding the $43$ aborted rows, the multiset period and the
      set-based period agree on every single one of the remaining
      $50{,}968$ rows except the row exhibited above.
\item The timeout cap of $120$\,s is heuristic; lifting it (or doubling
      \texttt{MAX\_PERIOD\_ARRAY\_SIZE}) might convert some of the $43$
      aborts into successful computations.
\end{itemize}

\end{document}
